\documentclass[11pt]{article}
\usepackage[margin=24mm]{geometry}
\usepackage{amsmath,amssymb,amsthm,mathtools,hyperref}
\newtheorem{theorem}{Theorem}
\newtheorem{lemma}[theorem]{Lemma}
\allowdisplaybreaks
\begin{document}
\title{A fixed testing family on a bounded C1 source interval}
\author{}\date{23 October 2026}\maketitle

\paragraph{The source and the complete parent.}
For positive $x,\sigma,z$ and $0<b,h<1$ put
\begin{gather*}
 r=bz<1,\quad Q=1-b,\quad\theta=1-h,\quad
 e=xh/\sigma,\quad E_s=1-e,\quad K_s=1+x,\quad R_s=1+xr,\\
 M_s=R_s-xz\theta,\quad J_s=1+xh-\sigma(1-r),\\
 P_s=\sigma+K_s+\sigma z(1+bx+2xh)
                          +\sigma bx^2hz^2+x^2h\theta z,\\
 \pi=\frac{z(1+xh)E_s+(1-r)(1+xhz)}{2P_s},\qquad
 l=\frac{(1+\sigma r)\{eR_s(1+\sigma z)+M_sE_s\}}{2P_s},\\
 \kappa_s=\frac{K_s(1+\sigma r)(1+\sigma z)}{2M_sP_s}.
\end{gather*}
The strict gaps are $E_s,M_s,J_s>0$. The coefficients are
\begin{gather*}
 m=\kappa_s x^2z h/\sigma,\quad
 C=\frac{K_sz(1+\sigma r)E_s}{2(1-r)P_s},\quad
 A=\frac\pi{bz},\quad \alpha=\frac l{Qz},\\
 D=\frac{K_s(1+\sigma z)(1+\sigma r)-J_sP_s}{2J_sP_s},\quad
 \beta=\frac{K_s(1+\sigma r)(1+xhz)-\theta P_s}{2\theta P_s}.
\end{gather*}
Require $D,\beta\ge0$. With $D_u=1+u^2$, $D_v=1+v^2$, write the
complete quarter-parent as
\begin{align*}
 \mathcal F={}&m(vW-X+uvY-uZ)^2+CD_v(uW-Y)^2+AD_uD_vY^2\\
 &+(D+\tfrac12)(uW+uvX-Y-vZ)^2
 +(\alpha+\tfrac12)D_u(vY-Z)^2\\
 &+(\beta+\tfrac12)D_v(X+uZ)^2-\tfrac12D_uD_vZ^2.
\end{align*}
The physical normalization is the pure testing factor $(D_uD_v)^{-1}$.
Homogeneous polynomial values define both projective angle endpoints.

\begin{theorem}\label{thm:bounded-unit}
Fix rationals $0<\ell\le U$. Define the following rational constants:
\begin{gather*}
 R=\max(1,1/\ell),\quad
 \varepsilon_R=2^{-40}(1+R)^{-8},\quad
 \delta_R=2^{-20}(1+R)^{-2},\\
 \mu=\frac1{4\max(U,1)},\quad
 M=\lceil96\max(U,1)\rceil,\quad\tau=\delta_R^2/2,\\
 \mathcal B=U^2(4+2M^2+M^4)+U(2+2M^2+M^4)+1+2M^2,\\
 k=\frac{\min(\ell,1)\tau^2}{8\mathcal B},\quad K=3M^2,\quad
 \varepsilon=\min\left(2^{-78},\varepsilon_R,
                                  \frac{k}{50(k+K)}\right).
\end{gather*}
Suppose the above native source satisfies
\[
 \ell\le L=1/z\le U,\qquad
 0<x,\sigma\le\frac\varepsilon{1+z}.
\]
There is one finite source-independent family of semialgebraic testing
functions whose source Gram assignments express the physical normalization of the complete parent
as a separated sum of squares at every pair of projective angles and
every state. This statement includes $D=0$, $\beta=0$, every allowed
source-opening rate, and all testing and state-coordinate zeros.
Its inner-corner input is the bounded-$R$ admission theorem on
$\sqrt{|u|},|w|\le\delta_R$, $w=1/v$, for native sources
$x,\sigma\le\varepsilon_R$, $0<z\le R$.
That theorem supplies fixed families of at most $220$ and $157$
functions in the two signs of its source angular coefficient.
\end{theorem}

All functions below are fixed after $\ell,U$ are chosen. They do not depend
on $L,r,h$, any event coefficient, or a source residue.

\paragraph{Quantitative native admission.}
Put $H_0=1+zQ$, $B_s=P_s+xz(J_s-\sigma)$ and
\[
 \mathcal H=R_sH_0+z\{(1-r)xh(2+xr)
                      -e[2+x(1+h+r\theta)]+K_se^2\}.
\]
Direct substitution in the displayed source formulas gives
\[
 2\theta P_s\beta=hB_s-\sigma H_0,\qquad
 2P_s(J_sD+x\theta\beta)=\sigma^2\mathcal H,
\]
and hence
\[
 eB_s=xH_0+\sigma\mathcal H-2P_sJ_sD/\sigma.
\]
Writing $T=1+z$ and using $Tx,T\sigma\le\varepsilon$,
$e<1$ and $r<1$ gives
\[
 1\le P_s\le1+4\varepsilon,\quad |M_s-1|\le\varepsilon,
 \quad B_s\ge1,\quad\mathcal H\le4T,
 \quad e\le5\varepsilon.
\]
For
\[
 C_0=\frac1{2L(1-r)},\quad
 A_0=\frac{1+L(1-r)}{2Lr},\quad
 a_0=\frac L{2(1-Lr)},
\]
the exact source ratios are
\begin{gather*}
 \frac C{C_0}=\frac{K_s(1+\sigma r)(1-e)}{P_s},\qquad
 \frac A{A_0}=\frac{z(1+xh)(1-e)+(1-r)(1+xhz)}{P_sH_0},\\
 \frac\alpha{a_0}=
 \frac{(1+\sigma r)\{eR_s(1+\sigma z)+M_s(1-e)\}}{P_s}.
\end{gather*}
The first two ratios lie between $1-9\varepsilon$ and
$1+3\varepsilon$; the third lies between $1-5\varepsilon$
and $1+5\varepsilon$. For the second numerator use
$H_0-ze+xhz(2-r-e)$ and $H_0-z=1-r>0$.
For the third use its convex combination of $M_s$ and
$R_s(1+\sigma z)$. In particular
\begin{equation}\label{eq:source-admit}
 d_s:=|C/C_0-1|+|A/A_0-1|+|\alpha/a_0-1|
 \le25\varepsilon\le\frac{k}{2(k+K)},\qquad C,A\ge\mu.
\end{equation}
No bound on $h/\sigma$ and no inverse of $D$ or $\beta$ was used.

\begin{lemma}[Both Bob charts at the Alice endpoint]\label{lem:endpoint}
If $C,A\ge\mu$ with $0<\mu\le1/4$, one fixed algebraic testing family
suffices whenever $0\le p=1/u\le\mu/24$, for every Bob angle.
The other nonnegative event coefficients are unrestricted.
\end{lemma}
\begin{proof}
On $|v|\le1$ use raw roots $(W,Y,B,R)$ with
$B=pX+Z$, $R=vX$. The complete homogeneous core is
\[
 K_p=B^2-2pvBW-2vBY+(1+p^2)R^2+2(1+p^2)RW
                       +(W-pY)^2+(1+p^2)v^2Y^2.
\]
Select only the fixed coefficients $C=A=\mu$. Their quarter-parent
\[
 Q_p=\mu(1+v^2)(W-pY)^2
          +\mu(1+p^2)(1+v^2)Y^2+K_p/2
\]
has the testing-only matrix
\[
 \begin{pmatrix}
 \mu D_v+1/2&-p(\mu D_v+1/2)&-pv/2&(1+p^2)/2\\
 -p(\mu D_v+1/2)&\mu D_v(1+2p^2)+(p^2+(1+p^2)v^2)/2&-v/2&0\\
 -pv/2&-v/2&1/2&0\\
 (1+p^2)/2&0&0&(1+p^2)/2
 \end{pmatrix}.
\]
At $p=0$ its form is
\[
 \mu D_v(W^2+Y^2)+\tfrac12(W+R)^2+\tfrac12(B-vY)^2.
\]
The raw norm is at most three times the squared norm of
$(W,Y,W+R,B-vY)$, so this matrix is at least $(\mu/3)I_4$.
The difference from its $p=0$ matrix has absolute row sums at most
$3p$ for $|v|\le1$, $p\le1$, $\mu\le1/4$.

On $|w|=|1/v|\le1$ use raw roots $(W,Y,X,H)$, $H=wZ$.
The complete twice-homogeneous core is
\begin{align*}
 K_{pw}={}&H^2-2pHW+2pwHX-2HY+w^2W^2+2wWX-2pw^2WY\\
 &+(1+p^2+p^2w^2)(X^2+Y^2)-2pwXY.
\end{align*}
The corresponding selected quarter-parent is
$Q_{pw}=\mu D_w(W-pY)^2+\mu(1+p^2)D_wY^2+K_{pw}/2$.
Its testing-only matrix is
\[
 \begin{pmatrix}
 \mu D_w+w^2/2&-p(\mu D_w+w^2/2)&w/2&-p/2\\
 -p(\mu D_w+w^2/2)&\mu D_w(1+2p^2)+(1+p^2D_w)/2&-pw/2&-1/2\\
 w/2&-pw/2&(1+p^2D_w)/2&pw/2\\
 -p/2&-1/2&pw/2&1/2
 \end{pmatrix}.
\]
At $p=0$ its form is
\[
 \mu D_w(W^2+Y^2)+\tfrac12(X+wW)^2+\tfrac12(H-Y)^2
                  \succeq(\mu/3)I_4.
\]
Its finite-$p$ correction has absolute row sums at most $4p$.
Thus both displayed matrices are at least $(\mu/6)I_4$ when
$p\le\mu/24$. Factoring them takes place entirely in the testing
field: they contain only fixed $\mu$ and testing variables.
Their four Cholesky roots in each chart are therefore fixed algebraic
testing functions. Every leading principal minor is strictly positive,
so the usual finite Cholesky formulas define them at all chart points.

Restore the six original nonnegative contributions with source
coefficients $C-\mu,A-\mu,\alpha,m,D,\beta$.
On the $v$ chart their fixed testing roots are
\begin{gather*}
 \sqrt{D_v}(W-pY),\quad\sqrt{(1+p^2)D_v}\,Y,
 \quad\sqrt{1+p^2}(vY-Z),\\
 pvW-pX+vY-Z,\quad W+vX-pY-pvZ,
 \quad\sqrt{D_v}(pX+Z).
\end{gather*}
On the $w$ chart they are
\begin{gather*}
 \sqrt{D_w}(W-pY),\quad\sqrt{(1+p^2)D_w}\,Y,
 \quad\sqrt{1+p^2}(Y-wZ),\\
 pW-pwX+Y-wZ,\quad wW+X-pwY-pZ,
 \quad\sqrt{D_w}(pX+Z).
\end{gather*}
This is at most ten fixed functions per chart, including their four
testing Cholesky roots. All original terms and literal endpoints
are included. The physical normalizers are pure testing.
\end{proof}

\paragraph{A finite annulus for all remaining bounded Alice angles.}
Put $c=(1+v^2)^{-1/2}$, $s=vc$, $P=uW$, $T=cZ$.
The selected compact-Bob quarter-parent is
\[
 S_0=C_0(P-Y)^2+A_0(1+u^2)Y^2
                       +a_0(1+u^2)(sY-T)^2+K_B/2,
\]
where
\[
 K_B=X^2+c^2(P-Y)^2+s^2Y^2-2sPT
 +2ucX(sP-sY+T)+u^2s^2X^2+u^2(sY-T)^2.
\]
For its raw Gram $G_0$ in $(P,X,Y,T)$, direct cofactors give
\[
 \operatorname{tr}G_0^{-1}
 =\frac{2(L^2N_2+LN_1+N_0)}
 {(1+t)\{Lq+t(1-q)\}\{1+L(1-r)\}},\quad t=u^2,\ q=c^2,
\]
with
\begin{align*}
 N_2&=1+r^2+(1-r)q\{1+r(1-q)\}
 +(1-r)(2-qr)t+(1-r)^2(1-q)t^2,\\
 N_1&=2(1-r)q+\{2-rq(3-q)\}t+(1-r)(1-q)t^2,\\
 N_0&=1+(2-q)t.
\end{align*}
The strict $T=0$ block and positive determinant establish
$G_0\succ0$ unless $u=c=0$. For $|u|\le M$,
\[
 N_2\le4+2M^2+M^4,\quad N_1\le2+2M^2+M^4,
 \quad N_0\le1+2M^2.
\]
Since $L\ge\ell$,
\[
 (1+t)\{Lq+t(1-q)\}\ge\min(\ell,1)(q+t).
\]
Consequently $G_0\succeq\min(\ell,1)\tau^2I/(2\mathcal B)$ on
$u^2+c^2\ge\tau^2$, and is at least $kI$ on the larger annulus
$u^2+c^2\ge\tau^2/2$. The complete core matrices, both in raw
coordinates and in
$e=(P-Y,Y,sY-T,X)$, have absolute row sums at most $K=3M^2$.
Using the common source factor $1-d_s$ and retaining all three event
excesses, \eqref{eq:source-admit} therefore gives an actual raw
reserve $k/2$ on that larger annulus.

Here is an explicit finite construction in separated fields. Define
\[
 N=K+1+24K^2/k,\qquad h_* =\min(1,k/24),\qquad B_*=N+1.
\]
Clip $C,A,\alpha$ at the fixed real $N$ and retain their excesses.
Both the event-coordinate map and its inverse have squared Frobenius
norm at most six. Thus the unmodified event-coordinate Gram is at
least $k/12$. Its core has norm at most $K$.
Splitting coordinates according to their source coefficient being
above $N$, the unchanged small block is at least $k/12$, the clipped
large block is at least $N-K$, and the off-diagonal norm is at most $K$.
The exact inequality
\[
 2Kab\le(24K^2/k)a^2+(k/24)b^2
\]
proves a clipped reserve $h_*I$.

The clipped event-coordinate Gram is affine in the seven testing
monomials
\[
 \eta=(u^2,c^2,s^2,s,ucs,uc,u^2s^2).
\]
Its coefficient matrices are exactly the following, with
$C',A',a'$ denoting the three clipped source weights:
\begin{gather*}
 G_{\rm const}=\operatorname{diag}(C',A',a',1/2),\qquad
 A_1=\operatorname{diag}(0,A',a'+1/2,0),\quad A_2=E_{11}/2,\\
 A_3=-(E_{12}+E_{21}+E_{22})/2,\quad
 A_4=(E_{13}+E_{31}+E_{23}+E_{32})/2,\\
 A_5=(E_{14}+E_{41}+E_{24}+E_{42})/2,\quad
 A_6=-(E_{34}+E_{43})/2,\quad A_7=E_{44}/2.
\end{gather*}
Each has norm at most $B_*$.
Choose the fixed integer and grid
\[
 n=\left\lceil\max\left(\frac{8M^2}{\tau^2},
                         \frac{56M^2B_*}{h_*}\right)\right\rceil,
 \quad \delta_g=4M^2/n,
\]
\[
 a_{jk}=\left(j/n,\frac{n^2-k^2}{n^2+k^2},
                         \frac{2kn}{n^2+k^2}\right),
 \quad -Mn\le j\le Mn,\quad -n\le k\le n.
\]
All centers are rational. The circle parametrization
$c=(1-t^2)/(1+t^2)$, $s=2t/(1+t^2)$, $-1\le t\le1$,
has both coordinate derivatives bounded by two. Thus the nearest
parameter grid point gives errors at most $1/n$ in $c,s$.
Retain centers with $u_j^2+c_k^2\ge\tau^2/2$.
Charts within $1/n$ of a center in each coordinate cover the target
annulus. On each chart $b_i=\eta_i-\eta_i(a)$ satisfies
$|b_i|\le\delta_g$: the derivative bounds for $u^2s^2$ give
$2M+2M^2\le4M^2$, and those of the other six monomials are smaller.
The exact identity
\[
 G'=G'(a)-7\delta_gB_*I
 +\tfrac12\sum_{i=1}^7
 \{(B_*I+A_i)(\delta_g+b_i)+(B_*I-A_i)(\delta_g-b_i)\}
\]
has positive source matrices, its first matrix being at least
$h_*I/2$. It uses only the fixed testing functions
$e_j$ and $e_j\sqrt{\delta_g\pm b_i}$.
Restore the fixed roots $\sqrt{1+u^2}Y$,
$\sqrt{1+u^2}(sY-T)$ and the three full $m,D,\beta$ events.
At most $65$ functions per chart express the complete parent.
No selected source/test matrix was square-rooted in a mixed field.

\paragraph{Completing the testing cover.}
When $|u|\ge M$, Lemma~\ref{lem:endpoint} applies because
$1/M\le\mu/24$. For $|u|\le M$ outside the displayed annulus,
$|u|,c<\tau$. Thus
\[
 |u|<\delta_R^2/2,\qquad
 |w|=c/\sqrt{1-c^2}<2\tau=\delta_R^2\le\delta_R.
\]
These inequalities imply the bounded-$R$ inner testing conditions
$\sqrt{|u|},|w|\le\delta_R$. Its source box contains the present box,
since $z\le1/\ell\le R$ and $x,\sigma\le\varepsilon\le\varepsilon_R$.

These are finitely many fixed testing charts. To obtain one family
on the whole testing domain, order the charts and assign a point to
the first chart containing it. Multiply each chart function by the
indicator of its resulting disjoint semialgebraic piece and extend it
by zero elsewhere. The identities then sum literally; all source
Gram assignments remain source-only. This introduces no additional
functions per chart. The indicators and radicals have fixed
semialgebraic graphs. Before this summation divide the compact-Bob
functions by $\sqrt{D_u}$, the endpoint $v$-chart functions by
$\sqrt{(1+p^2)D_v}$, and the endpoint $w$-chart functions by
$\sqrt{(1+p^2)D_w}$. Use the corresponding physical normalizer for
the inner family, with its constant factor $1/2$ if its target is
the full parent $4\mathcal F$. All charts then represent the same
quadratic function $\mathcal F/(D_uD_v)$, including its homogeneous
projective values. The exact sign gauges are
\[
 (u,W,X,Y,Z)\longmapsto(-u,W,X,-Y,-Z),\qquad
 (v,W,X,Y,Z)\longmapsto(-v,W,-X,Y,-Z).
\]
They preserve every square term and the signed complete core, so
negative $p$ requires no additional source statement.
Real/imaginary splitting handles complex
states. This proves Theorem~\ref{thm:bounded-unit}.

\paragraph{Scope.}
The fixed testing family is obtained for each specified interval
$[\ell,U]$. It is not uniform over all positive real residues of $L$
without fixing such an interval, and does not treat $z=0$ or other
source charts. The bounded-$R$ inner functions depend on the fixed $R$;
this theorem does not identify them with the evaluated order-$731$
bounded-$z$ construction.
When $\ell\ge1$, one may instead use the published local669 theorem,
choose $\tau=2^{-25}$, and omit $\varepsilon_R$ from the minimum.
Then the same proof requires no bounded-$R$ input.
This is a C1 source-domain result, not a global semidefinite lift
of the quantum set.
\end{document}
